\section*{AI-use statement}

This paper was drafted by a crew of AI research-assistant agents
(Free-Research-Crew, protocol v0.2) under the direction of the human author, Bryce Wang.
He chose the topic from a proposal and approved the research question (2026-11-21, 05:35
UTC), the design (2026-11-21, 11:41 UTC) and the results (2026-11-21, 14:12 UTC), each in
chat after reading the Reviewer's verdict. The manuscript itself is subject to the
submission gate. He takes responsibility for the paper's content.

% Counts are macros from process.tex (`crew.py numbers`); who did what is taken from
% log.jsonl, reviews/ and briefs/.
\begin{itemize}
  \item \textbf{Roles and models.} Every role (PI, Scout, Designer, Data RA, Econometrician,
        Figures RA, Writer, Reviewer and Replicator) ran on Claude Opus 5.5
        (\texttt{claude-opus-5-5}) in Claude Code, each role as a separate agent. The PI
        proposed the topic, drafted the question and hypotheses, wrote the contribution
        statement, recorded deviations D5--D8 and corrected threat T14 after the results
        gate; the Scout found and verified the sources; the Designer wrote the design and
        deviations D1--D4; the Data RA fetched and prepared the data; the Econometrician
        wrote and tested the estimation code, executed the analysis runs, including the
        post hoc and descriptive runs, and wrote the result claims; the Figures RA wrote the figure scripts
        and drew the figures; the Writer drafted this paper from the claims; the Replicator
        ran the clean rerun. The PI's brief to each role is archived in \texttt{briefs/}.
  \item \textbf{Review.} The Reviewer ran \CrewReviewRounds{} rounds on
        \CrewReviewerModels{} (\CrewReviewIndependence{}: a new context in every round, the
        same model family as the other roles): \CrewRoundsQuestion{} on the question,
        \CrewRoundsLiterature{} on the literature, \CrewRoundsDesign{} on the design,
        \CrewRoundsCode{} on the code, \CrewRoundsResults{} on the results and
        \CrewRoundsPaper{} on the paper. It raised \CrewFindingsCritical{} critical,
        \CrewFindingsMajor{} major and \CrewFindingsMinor{} minor distinct findings (a
        finding carried into a later round is counted once, at its latest severity). The
        last round of each completed stage passed. Four sets of changes were made after
        the last round of their stage and reached the human at the next gate without
        another round of that stage: the design's responses to the check-methods failures
        (the SDID critical-value floor, the two-part synthetic-control check and the
        adopted implementation choices), which the code review then checked against the
        code; two report-only additions to the design (findings D14 and D15); claim-text corrections after the third results round (findings R12 and R13); and
        the PI's correction of threat T14 after the results gate (the 2025 quarters are final), which changes no analysis.
  \item \textbf{Runs.} \CrewRuns{} runs in the ledger (\CrewRunsAnalysis{} analysis,
        \CrewRunsDataPrep{} data preparation); \CrewRunsSuperseded{} were replaced by a
        later run with the same id and archived (the first data fetch, replaced to drop two metrics that a clean rerun could not reproduce, and the
        first run of the method checks, whose failures led to the design changes described above; no run behind
        a result claim was replaced); \CrewRunsUnreviewed{} analysis runs were not covered
        by a passing code review, both after the fourth and last code-review round and both
        with the human author's consent in chat: the post hoc run (deviation D7), whose
        code was changed to add the leave-one-out check the Reviewer asked for, and the
        descriptive run behind the summary-statistics table (deviation D8), added at the
        paper stage; the paper review read both runs' outputs.
  \item \textbf{Replication.} The Replicator's clean rerun (\texttt{crew.py rerun}, in a
        fresh environment built from \texttt{code/requirements.txt}) has status
        \CrewRerunStatus{}: \CrewRerunReproduced{} of \CrewRerunTotal{} runs reproduced.
        Whether the work can be rerun from scratch is what this record states, not
        something the authoring roles assert.
  \item \textbf{Not blind.} The PI and the Scout knew the published estimates and a
        published descriptive index of California fast-food employment and weekly wages
        through 2026Q1 before the design was fixed; the primary specification was
        therefore chosen by a rule on data through 2019, and the analysing roles computed
        nothing on post-2023Q3 outcomes before the design gate.
  \item \textbf{Findings accepted over the Reviewer's objection.} None.
  \item \textbf{Evidence.} Every result number in the text, the tables and the figure
        captions is a macro generated from the metrics of a logged run. The numbers typed
        by hand are design facts fixed before any analysis run, published numbers quoted
        from the literature as the Scout's source notes record them, and process facts
        such as dates. The full workspace, including every brief, review and run record, is
        published with the paper.
\end{itemize}
